% from ega3-1-fr.tex
\oldpage[0\textsubscript{III}]{43}
\item[{\rm(ii)}] Soit $M_{\bullet\bullet\bullet}$ un tricomplexe de
$\mathcal C$ tel que pour chaque couple $(r,s)$,
$M_{rs,\bullet}$ soit une résolution de $K_{rs}$ et que
$L_nT(M_{ijk})=0$ pour tous les triplets $(i,j,k)$ et tout $n>0$. Alors on a
un isomorphisme
\[
L_\bullet T(K_{\bullet\bullet})
\simeq H_\bullet(T(M_{\bullet\bullet\bullet}));
\]
en effet, avec les notations de la démonstration de
(\hyperref[0.11.7.2-fr]{11.7.2}), $N_{i,\bullet}$ est une résolution de
$L_i$ telle que $L_nT(N_{ij})=0$ pour $n>0$ et tout couple $(i,j)$, et il
suffit d'appliquer (\hyperref[0.11.6.3-fr]{11.6.3}, (iii)).
\item[{\rm(iii)}] On généralise aussitôt les résultats de
(\hyperref[0.11.7.2-fr]{11.7.2}) aux multifoncteurs covariants; par exemple,
soient $\mathcal C'$ une seconde catégorie abélienne dans laquelle tout objet
est quotient d'un objet projectif, $K'_{\bullet\bullet}$ un bicomplexe de
$\mathcal C'$ dont les deux degrés sont limités inférieurement, et $T$ un
bifoncteur covariant additif de $\mathcal C\times\mathcal C'$ dans une
catégorie abélienne $\mathcal C''$. Si $L_\bullet$ et $L'_\bullet$ sont les
complexes simples associés à $K_{\bullet\bullet}$ et
$K'_{\bullet\bullet}$ respectivement, on définit l'hyperhomologie de $T$ par
rapport aux deux bicomplexes $K_{\bullet\bullet}$,
$K'_{\bullet\bullet}$ comme l'hyperhomologie
$L_\bullet T(L_\bullet,L'_\bullet)$; appliquant
(\hyperref[0.11.6.4-fr]{11.6.4}) et
(\hyperref[0.11.6.5-fr]{11.6.5}), on a de même que dans
(\hyperref[0.11.7.2-fr]{11.7.2}) six suites spectrales aboutissant à cette
hyperhomologie, et que nous laissons le soin d'écrire au lecteur.
\end{enumerate}
\end{env}

\subsection{Compléments sur la cohomologie des complexes simpliciaux}
\label{subsection:0.11.8-fr}

\begin{env}[11.8.1]
\label{0.11.8.1-fr}
Soient $A$ un ensemble fini, $\Sigma(A)$ l'ensemble des suites finies
$\sigma=(\alpha_0,\ldots,\alpha_h)$ d'éléments de $A$ («~simplexes~» de
$A$); on pose $|\sigma|=\{\alpha_0,\ldots,\alpha_h\}$; on rappelle que le
complexe de chaînes $\mathrm C_\bullet(A)$ est le groupe abélien libre
gradué engendré par les éléments de $\Sigma(A)$,
$(\alpha_0,\ldots,\alpha_h)$ étant de degré $h$, avec une différentielle
définie par
\[
d(\alpha_0,\ldots,\alpha_h)=
\sum_{j=0}^{h}(-1)^j
(\alpha_0,\ldots,\widehat{\alpha_j},\ldots,\alpha_h).
\]
Le sous-groupe $\mathrm D_\bullet(A)$ de $\mathrm C_\bullet(A)$ engendré par
les chaînes $\sigma=(\alpha_0,\ldots,\alpha_h)$ pour lesquelles deux des
$\alpha_i$ sont égaux, et par les chaînes
$\pi(\sigma)-\varepsilon_\pi\sigma$, où pour toute permutation $\pi$,
\[
\pi(\sigma)=(\alpha_{\pi(0)},\ldots,\alpha_{\pi(h)})
\]
et $\varepsilon_\pi$ est la signature de $\pi$, est un sous-complexe de
$\mathrm C_\bullet(A)$ dont les éléments sont dits \emph{chaînes
dégénérées}; on pose
$\mathrm L_\bullet(A)=\mathrm C_\bullet(A)/\mathrm D_\bullet(A)$, et on a un
homomorphisme naturel de complexes
\[
p:\mathrm C_\bullet(A)\longrightarrow\mathrm L_\bullet(A).
\]
On définit d'autre part un homomorphisme de complexes
\[
j:\mathrm L_\bullet(A)\longrightarrow\mathrm C_\bullet(A)
\]
de la façon suivante : on ordonne totalement $A$; à la classe
mod.~$\mathrm D_\bullet(A)$ d'un simplexe
$\sigma=(\alpha_0,\ldots,\alpha_h)$, on fait correspondre $0$ si deux des
$\alpha_i$ sont égaux, et la suite $(\beta_0,\ldots,\beta_h)$ des $\alpha_i$
rangés dans l'ordre croissant dans le cas contraire. Il est clair que $p\circ
j$ est l'identité de $\mathrm L_\bullet(A)$.
\end{env}

\begin{env}[11.8.2]
\label{0.11.8.2-fr}
Soit $B$ un second ensemble fini; si $d'$, $d''$ sont les différentielles de
$\mathrm C_\bullet(A)$ et $\mathrm C_\bullet(B)$, le complexe produit
tensoriel $\mathrm C_\bullet(A)\otimes\mathrm C_\bullet(B)$ peut être
considéré comme le groupe abélien libre engendré par les éléments de
$\Sigma(A)\times\Sigma(B)$, avec la différentielle
\[
d(\sigma,\tau)=(d'\sigma,\tau)+(-1)^{h+1}(\sigma,d''\tau)
\quad\text{si }\operatorname{Card}|\sigma|=h+1.
\]
Les homomorphismes naturels
$\mathrm C_\bullet(A)\to\mathrm L_\bullet(A)$,
$\mathrm C_\bullet(B)\to\mathrm L_\bullet(B)$ définissent un homomorphisme
\[
p:\mathrm C_\bullet(A)\otimes\mathrm C_\bullet(B)
\longrightarrow\mathrm L_\bullet(A)\otimes\mathrm L_\bullet(B),
\]
ce dernier produit tensoriel étant isomorphe à
\[
\frac{\mathrm C_\bullet(A)\otimes\mathrm C_\bullet(B)}
{\mathrm C_\bullet(A)\otimes\mathrm D_\bullet(B)
+\mathrm D_\bullet(A)\otimes\mathrm C_\bullet(B)}.
\]
De même, à l'aide des homomorphismes
$\mathrm L_\bullet(A)\to\mathrm C_\bullet(A)$ et
$\mathrm L_\bullet(B)\to\mathrm C_\bullet(B)$ définis dans
(\hyperref[0.11.8.1-fr]{11.8.1}) (à l'aide d'ordres totaux sur $A$ et $B$),
on définit un homomorphisme
\[
j:\mathrm L_\bullet(A)\otimes\mathrm L_\bullet(B)
\longrightarrow\mathrm C_\bullet(A)\otimes\mathrm C_\bullet(B)
\]
tel que $p\circ j$ soit l'identité.
\end{env}

